\begin{frame}<1-4>[t]{新知探索}
\begin{campuscanvas}
% Source slide 13, objects 25; visible 2-4
\only<2-4|handout:1>{\node[anchor=north west,inner sep=0,text=black] at (70.000,150.000) {\begin{minipage}[t]{142.500mm}\raggedright\fontsize{16}{24.00}\selectfont \textbf{5. 表示集合的方法：列举法}\end{minipage}};}
% Source slide 13, objects 24; visible 3-4
\only<3-4|handout:1>{\node[anchor=north west,inner sep=0,text=black] at (70.000,245.000) {\begin{minipage}[t]{142.500mm}\raggedright\fontsize{12}{18.00}\selectfont 表示一个集合，就是把它有哪些元素交代清楚。生活中常见的方法，是\red{把集合中的元素一一列举出来}。\par 饭馆里的菜单，计算机里的文件夹，各种委员会名单，都是这样做的。这种表示法叫作\red{列举法}。\end{minipage}};}
% Source slide 13, objects 5; visible 4
\only<4|handout:1>{\node[anchor=north west,inner sep=0,text=black] at (70.000,475.000) {\begin{minipage}[t]{142.500mm}\raggedright\fontsize{12}{18.00}\selectfont 数学里用列举法表示集合，常用的格式是在一个大括号里写出每个元素的名字，相邻的名字用逗号分隔。\par 例如，小于 $10$ 的正偶数组成的集合，用列举法可以表示为\[\{2,4,6,8\}\quad\text{或}\quad\{8,2,6,4\}\quad\text{等。}\]\end{minipage}};}
\end{campuscanvas}
\end{frame}
